\chapter{Calabi--Yau Structures and Moduli of Objects}
\label{cygeom:chap:moduli}

An extension varies, has automorphisms, and can meet obstructions to further
variation. Calabi--Yau duality relates these three features. On a threefold it
pairs first-order deformations with their obstruction spaces and pairs
automorphisms with degree-three endomorphisms. These pairings determine the
shifted tangent geometry and determinant lines that enter extension
correspondences. Their construction retains both the underlying complexes
and the coherence that makes the pairings vary over a moduli space.

Throughout, $k$ is a field of characteristic zero. Complexes are
cohomologically graded, with $V[r]^i=V^{i+r}$ and
$d_{V[r]}=(-1)^r d_V$. The dual of a complex is
$V^*=\operatorname{RHom}_k(V,k)$. A perfect $k$-complex has finite-dimensional
total cohomology. All tensor products and Hom complexes are derived unless
their displayed free models make this unnecessary.

\section{The deformation spaces of a point}
\label{cygeom:sec:point}

Let $A=k[x_1,x_2,x_3]$ and let $E=A/(x_1,x_2,x_3)$ be the module of the
origin. Write $W=k e_1\oplus k e_2\oplus k e_3$. The Koszul resolution of
$E$ is
\begin{equation}\label{cygeom:eq:koszul}
 K^{-p}=A\otimes\bigwedge^p W,\qquad
 d(e_{i_1}\wedge\cdots\wedge e_{i_p})
 =\sum_{a=1}^p(-1)^{a-1}x_{i_a}
 e_{i_1}\wedge\cdots\widehat e_{i_a}\cdots\wedge e_{i_p}.
\end{equation}

Here each $e_i$ has cohomological degree $-1$. The sequence
$x_1,x_2,x_3$ is regular: after quotienting by its first $i-1$ terms,
$x_i$ is still a non-zero-divisor. Successively tensoring the corresponding
two-term resolutions therefore proves that $K\to E$ is a resolution.
After applying $\operatorname{Hom}_A(-,E)$, its differential is zero.

\begin{proposition}\label{cygeom:prop:point-ext}
Let $\xi_1,\xi_2,\xi_3$ be the basis dual to $e_1,e_2,e_3$.
The Yoneda algebra of the origin is
\[
 \operatorname{Ext}^*_A(E,E)
   =\bigwedge^*(k\xi_1\oplus k\xi_2\oplus k\xi_3),
 \qquad |\xi_i|=1.
\]
The functional $\tau(\xi_1\xi_2\xi_3)=1$, zero in other degrees,
gives a perfect pairing
$\operatorname{Ext}^p_A(E,E)\otimes
 \operatorname{Ext}^{3-p}_A(E,E)\to k$ by multiplication.
\end{proposition}

\begin{proof}
The resolution gives the stated graded vector space. To compute its
product, let $\iota_i$ be contraction with $e_i^*$ on the exterior factor
of $K$. It is an $A$-linear endomorphism of degree one.
The relations $\iota_i\iota_j+\iota_j\iota_i=0$ and
$d=\sum_i x_i\iota_i$ show that each $\iota_i$ is closed in
$\operatorname{End}_A(K)$.

Composition with $K\to E$ identifies their
products with $(-1)^{p(p-1)/2}$ times the dual exterior basis in degree $p$ of
$\operatorname{Hom}_A(K,E)$. This composition is a quasi-isomorphism
because $K$ is a bounded free resolution. It computes the Yoneda product
and proves the exterior presentation. A basis monomial pairs only with
its complementary monomial, with coefficient $1$ or $-1$.
\end{proof}

Thus the four dimensions are $1,3,3,1$. The three parameters moving the
point occur in degree one. Their three duals occur in degree two.
Degree zero and degree three supply the remaining pair. This computation
already distinguishes a three-dimensional space of ordinary points from
the derived moduli of their modules. We next identify the categorical
structure responsible for its duality.

\section{Two finiteness conditions}
\label{cygeom:sec:finiteness}

A differential graded category has complexes of morphisms and composition
compatible with their differentials. We use small $k$-linear dg categories,
and pass to perfect modules when shifts, cones and direct summands are
needed. For a dg algebra $A$, the category $\operatorname{Perf}(A)$ consists
of the direct summands of finite iterated cones of shifts of the free
right module $A$. A \emph{generator} is an object whose shifts, cones and
direct summands generate the category in this sense.

These operations make an algebra presentation useful without making it
intrinsic. If $G$ generates a dg category and
$A=\operatorname{RHom}(G,G)$, the functor
$\operatorname{RHom}(G,-)$ identifies its perfect completion with
$\operatorname{Perf}(A)$. The identification of morphism complexes
holds on $G$ and extends under cones and summands in both variables.
Its image is precisely the subcategory generated by $A$. An equivalence of
perfect module categories is called a \emph{Morita equivalence}.

\begin{definition}\label{cygeom:def:finiteness}
A dg category $\mathcal C$ is \emph{proper} if
$\operatorname{RHom}_{\mathcal C}(E,F)$ is perfect over $k$ for every
$E,F$. It is \emph{smooth} if its diagonal bimodule
$(E,F)\mapsto\operatorname{RHom}_{\mathcal C}(E,F)$ is perfect over
$\mathcal C^{\mathrm{op}}\otimes\mathcal C$.
\end{definition}

Properness concerns the size of each space of morphisms. Smoothness
concerns a finite resolution of the identity functor. For an algebra
$A$, put $A^e=A\otimes A^{\mathrm{op}}$, with the graded opposite
and the Koszul tensor signs. Smoothness of $\operatorname{Perf}(A)$
means that the diagonal $A$ is perfect as an $A^e$-module. Properness
means that $A$ is perfect over $k$.

The polynomial algebra $A=k[x_1,\ldots,x_d]$ is smooth: the Koszul
complex on $x_i\otimes1-1\otimes x_i$ resolves its diagonal over $A^e$.
For $d>0$ it is not proper, since
$\operatorname{Hom}_A(A,A)=A$ is infinite-dimensional. The converse
independence also has a small example.

A Serre functor on a proper category is an autoequivalence $S$ with
natural isomorphisms
$\operatorname{RHom}(E,F)^*\simeq\operatorname{RHom}(F,SE)$.

\begin{proposition}\label{cygeom:prop:dual-numbers}
For $D=k[\epsilon]/(\epsilon^2)$, the category
$\operatorname{Perf}(D)$ is proper and has the identity as its Serre
functor, but is not smooth.
\end{proposition}

\begin{proof}
The functional $\tau_D(a+b\epsilon)=b$ has multiplication pairing
$\left(\begin{smallmatrix}0&1\\1&0\end{smallmatrix}\right)$ in the basis
$1,\epsilon$. It identifies $D$ with its linear dual as a bimodule.
For bounded complexes of finite projective modules, evaluation gives
\[
 \operatorname{RHom}_D(P,Q)^*
   \simeq\operatorname{RHom}_D(Q,P\otimes_D D^*)
   \simeq\operatorname{RHom}_D(Q,P).
\]
The first identification can be checked on free modules and then extended
under cones and summands. All these Hom complexes are perfect over $k$.
This proves properness and the asserted Serre duality.

For smoothness, set $B=D^e=k[x,y]/(x^2,y^2)$. The diagonal is the
quotient sending $x,y$ to $\epsilon$, and has the free resolution
\begin{equation}\label{cygeom:eq:periodic}
 \cdots\xrightarrow{x-y}B\xrightarrow{x+y}B
 \xrightarrow{x-y}B\longrightarrow D\longrightarrow0.
\end{equation}
Both consecutive products vanish. In the basis $1,x,y,xy$,
\[
 \ker(x-y)=\langle x+y,xy\rangle=\operatorname{im}(x+y),\qquad
 \ker(x+y)=\langle x-y,xy\rangle=\operatorname{im}(x-y).
\]
The kernel of $B\to D$ is also $\langle x-y,xy\rangle$, proving
exactness.

Tensoring with $B/(x,y)=k$ kills every differential, so
$\operatorname{Tor}^B_i(D,k)=k$ for every $i\geq0$.
A perfect diagonal would give zero Tor in all sufficiently high degrees.
\end{proof}

Smoothness and properness together ensure its existence on the perfect
category. The proposition shows that smoothness is not necessary for
existence. The distinction matters when a Calabi--Yau condition is stated
using a Serre functor.

\section{A volume class and a trace}
\label{cygeom:sec:structures}

There are two duals of the diagonal. In an algebra presentation they are
\[
 A^!=\operatorname{RHom}_{A^e}(A,A^e),\qquad
 A^*=\operatorname{RHom}_k(A,k).
\]
In the first expression the resulting right $A^e$-module is regarded as
a bimodule by exchanging the two factors with the graded symmetry.
When both finiteness conditions hold, tensoring with $A^*$ is the Serre functor and
tensoring with $A^!$ is its inverse
\cite{cygeom:BD-relative}*{Proposition~2.3 and Example~2.12}.

To state the coherent duality datum, let $C(A)$ denote the Hochschild
chain complex. We write $C_j(A)$ for its cohomological degree $-j$.
Its Hochschild differential $b$ lowers $j$ by one. The Connes operator
$B$ raises $j$ by one and satisfies
\[
 b^2=B^2=bB+Bb=0.
\]
A complex with these two operators is a \emph{mixed complex}.
With a formal variable $u$ of homological degree $-2$, define the
negative cyclic and cyclic complexes by
\begin{equation}\label{cygeom:eq:cyclic-complexes}
 CC^-(A)=(C(A)[[u]],b+uB),\qquad
 CC(A)=(C(A)((u))/uC(A)[[u]],b+uB).
\end{equation}

The first uses products in each total degree. In the second, each
representative has only finitely many nonpositive powers of $u$.
These complexes model the homotopy fixed points and homotopy orbits,
respectively, of the circle action on Hochschild chains. Their maps below
use our cohomological shift convention.

\begin{definition}\label{cygeom:def:cy-structures}
A \emph{left Calabi--Yau structure of dimension $d$} on a smooth dg
category is a class $\eta:k[d]\to CC^-(\mathcal C)$ with the following
nondegeneracy property. Its underlying Hochschild class induces
an equivalence of bimodules
\[
 \mathcal C^!\xrightarrow{\sim}\mathcal C[-d].
\]
A \emph{right Calabi--Yau structure of dimension $d$} on a proper dg
category is a map
$\tau:CC(\mathcal C)\to k[-d]$ whose underlying Hochschild functional
induces an equivalence
\[
 \mathcal C[d]\xrightarrow{\sim}\mathcal C^*.
\]
These are the left and right structures of
\cite{cygeom:BD-relative}*{Definitions~3.2 and~3.5}.
\end{definition}

The maps have different variances: $\eta$ is a homology class, whereas
$\tau$ is a functional. Smoothness identifies Hochschild chains with
$\operatorname{RHom}_{A^e}(A^!,A)$, which explains the map induced by
$\eta$. Dualizing Hochschild chains identifies their functionals with
$\operatorname{RHom}_{A^e}(A,A^*)$, which explains the map induced by
$\tau$. If one retains only these nondegenerate Hochschild data, the
structure is called \emph{weak}. A cyclic lift is part of a left or
right structure as defined here.

On a smooth proper category the underlying left and right dualities
give the same Serre shift $S\simeq[d]$, since the two dualizing functors
are inverse. The coherent structures are related by the circle-equivariant
evaluation pairing. An object $E$ is \emph{proper} if
$\operatorname{RHom}(F,E)$ is perfect over $k$ for every perfect $F$.
A left structure induces a right structure on the full subcategory of
proper objects \cite{cygeom:BD-relative}*{Theorem~3.1}. This is the duality
used for objects with proper support in a nonproper space.

The negative cyclic condition is visible as a sequence of equations.
Writing a degree-$d$ cycle as
$\eta=\eta_0+u\eta_1+u^2\eta_2+\cdots$, where
$\eta_r\in C_{d+2r}(A)$, gives
\begin{equation}\label{cygeom:eq:lift}
 b\eta_0=0,\qquad b\eta_{r+1}=-B\eta_r\quad(r\geq0).
\end{equation}
Thus the first additional requirement is a primitive for $B\eta_0$.

For example, take a mixed complex with basis $e,f$ in homological degrees
$0,1$, with $b=0$ and $Be=f$. Its class $e$ has no negative cyclic
lift. Adjoining $g$ in degree two with $bg=-f$ and $Bg=0$ gives
the lift $e+ug$. This example is a statement about mixed complexes.
It makes no assertion that the two complexes are Hochschild complexes
of smooth proper categories.

\section{The polynomial volume class}
\label{cygeom:sec:volume}

For $A=k[x_1,\ldots,x_d]$, the left structure can be constructed in
coordinates. In normalized Hochschild chains write
$a_0[a_1|\cdots|a_n]$. A unit in any bracketed slot gives zero.
The degree-$d$ chain
\begin{equation}\label{cygeom:eq:volume-cycle}
 \eta_0=\sum_{\sigma\in S_d}\operatorname{sgn}(\sigma)
       1[x_{\sigma(1)}|\cdots|x_{\sigma(d)}]
\end{equation}
is a negative cyclic cycle with no higher $u$-terms.

\begin{proposition}\label{cygeom:prop:polynomial-cy}
The class \eqref{cygeom:eq:volume-cycle} defines a left Calabi--Yau
structure of dimension $d$ on $\operatorname{Perf}(A)$.
Its underlying bimodule equivalence is $A^!\simeq A[-d]$.
\end{proposition}

\begin{proof}
The adjacent multiplication terms in $b\eta_0$ cancel in pairs by
commutativity and antisymmetry. The two end terms cancel as well.
Every term of $B\eta_0$ has the initial unit of $\eta_0$ in a bracketed
slot, and therefore vanishes in the normalized complex. This proves
the negative cyclic equations.

Let $K_\Delta$ be the free bimodule Koszul resolution on
$x_1-y_1,\ldots,x_d-y_d$. The regular-sequence argument used for
\eqref{cygeom:eq:koszul} applies to this sequence too. Its dual
$\operatorname{Hom}_{A^e}(K_\Delta,A^e)$ has cohomology only in degree
$d$, where it is $A$.

The antisymmetrization map from this Koszul
resolution to the bar resolution sends its top generator to the sum
in \eqref{cygeom:eq:volume-cycle} after tensoring with the diagonal.
Indeed, its differential telescopes to the terms $x_i-y_i$. All
interior products cancel in adjacent transpositions. The induced map
from the dual Koszul complex sends the dual top generator to $1$ in
$A[-d]$. It is an isomorphism on its only cohomology group, proving
nondegeneracy and fixing the normalization.
\end{proof}

Write $\HH^p(A)=H^p\operatorname{RHom}_{A^e}(A,A)$ and
$\HH_q(A)=H^{-q}(A\otimes_{A^e}^{\mathbf L}A)$.
The same resolution computes
\[
 \HH^p(A)=A\otimes\bigwedge^p\langle\partial_1,\ldots,\partial_d\rangle,
 \qquad
 \HH_q(A)=A\otimes\bigwedge^q\langle dx_1,\ldots,dx_d\rangle.
\]
Cap product with $dx_1\wedge\cdots\wedge dx_d$ identifies
$\HH^p(A)$ with $\HH_{d-p}(A)$ by contraction. For example,
$\iota_{\partial_i}(dx_1\wedge\cdots\wedge dx_d)
=(-1)^{i-1}dx_1\wedge\cdots\widehat{dx_i}\cdots\wedge dx_d$.

This duality changes how the degrees are displayed. The graded
commutator of cochain insertions defines the Hochschild Gerstenhaber
bracket, of degree $-1$. Its value on a derivation and a function is
evaluation, so $[\partial_i,x_j]=\delta_{ij}$. Deforming an object instead
introduces a shift of its endomorphism complex.

\section{The tangent complex and its pairing}
\label{cygeom:sec:tangent}

For a connective commutative dg $k$-algebra $R$, a \emph{pseudo-perfect
family} over $R$ is a dg functor
$\mathcal C^{\mathrm{op}}\to\operatorname{Perf}(R)$.
The derived space of such functors and their equivalences defines the
moduli prestack $\mathcal M_{\mathcal C}(R)$. In an algebra presentation,
it is the moduli of right $A\otimes R$-modules whose underlying
$R$-complex is perfect. This condition concerns the family over the
base. For a smooth category its fibers are perfect modules over the
category as well.

We impose a finite-type hypothesis for the geometric statements below.
The perfect category of $\mathcal C$ is of \emph{finite type} if it is
compact in the Morita category of small idempotent-complete dg categories.
This means that mapping spaces out of it commute with filtered colimits.
Under this hypothesis $\mathcal M_{\mathcal C}$ is locally a derived
Artin stack of finite presentation and has a perfect cotangent complex
\cite{cygeom:BD-moduli}*{Remark~3.9}. Its cotangent complex
$\mathbb L_{\mathcal M}$ represents derivations into square-zero
extensions. Its dual is the tangent complex $\mathbb T_{\mathcal M}$.

The affine-three-space example satisfies the finite-type hypothesis.
One finite free dg algebra model has generators
$x_1,x_2,x_3$ in degree zero, $h_{12},h_{13},h_{23}$ in degree $-1$,
and $t$ in degree $-2$, with differential
\begin{equation}\label{cygeom:eq:finite-model}
 \begin{split}
 dx_i&=0,\qquad dh_{ij}=x_ix_j-x_jx_i,\\
 dt&=[x_1,h_{23}]-[x_2,h_{13}]+[x_3,h_{12}].
 \end{split}
\end{equation}
Here $[x,h]=xh-hx$, and $d$ extends by the graded Leibniz rule.
The Jacobi identity gives $d^2t=0$. The map to
$A=k[x_1,x_2,x_3]$ sends $h_{ij},t$ to zero and is a quasi-isomorphism.

For the augmentation $x_i\mapsto0$, the reduced bar construction of
$A$ has words in its augmentation ideal shifted by one. Its differential
multiplies adjacent entries, and its coproduct cuts words. The exterior
coalgebra on three degree-$-1$ primitive elements maps to it by
antisymmetrization. The Koszul
resolution of $k$ proves that this map induces an isomorphism on Tor,
hence is a quasi-isomorphism.

The free dg algebra
\eqref{cygeom:eq:finite-model} is its cobar: desuspend the reduced
coalgebra and extend the reduced coproduct as a differential.
Give $x_i,h_{ij},t$ weights $1,2,3$, respectively.
Cobar preserves this quasi-isomorphism, as the tensor-length filtration
is finite in each polynomial weight. Bar--cobar contraction then gives
the asserted map to $A$. There are only seven free generators, attached
in the displayed order. Maps out of this finite cell presentation commute
with filtered colimits, and so do maps out of its perfect category.

At an object $E$ there is a natural equivalence
\begin{equation}\label{cygeom:eq:tangent}
 \mathbb T_{\mathcal M,E}\simeq
 V_E[1],\qquad V_E=\operatorname{RHom}_{\mathcal C}(E,E)
 \quad\text{\cite{cygeom:BD-moduli}*{Proposition~3.3}}.
\end{equation}

The shift can be seen directly. On a cofibrant module representative,
a first-order change of differential is $d_E+\varepsilon a$, with
$|a|=1$ and $\varepsilon^2=0$. Squaring gives
$d_Ea+ad_E=0$. Conjugation by $1+\varepsilon h$, with $|h|=0$,
changes $a$ by a boundary. Thus degree-one endomorphisms describe
deformations, degree-zero endomorphisms describe their infinitesimal
automorphisms, and degree-two endomorphisms receive the obstruction to
lifting a first-order deformation. The derived statement
\eqref{cygeom:eq:tangent} retains this entire complex.

Suppose that $E$ carries the proper trace induced by a $d$-dimensional
Calabi--Yau structure. It gives a map
\[
 \tau_E:V_E\longrightarrow k[-d]
\]
whose cohomology pairing is perfect and satisfies
$\tau_E(xy)=(-1)^{|x||y|}\tau_E(yx)$.
When we compute with a cyclic chain model, this identity holds on that
model. Put $sx\in V_E[1]$
for the shift of a homogeneous element $x$. Thus $|sx|=|x|-1$.
The corresponding tangent pairing is
\begin{equation}\label{cygeom:eq:suspended-pairing}
 \omega_E(sx,sy)=(-1)^{|x|}\tau_E(xy)
       \;:\;V_E[1]\otimes V_E[1]\longrightarrow k[2-d].
\end{equation}

\begin{lemma}\label{cygeom:lem:shifted-pairing}
On cohomology, the pairing \eqref{cygeom:eq:suspended-pairing} is graded
alternating and nondegenerate. In a cyclic chain model it is an
alternating map of complexes.
\end{lemma}

\begin{proof}
Write $p=|x|$ and $q=|y|$. Cyclicity gives
\[
 -(-1)^{(p-1)(q-1)}\omega_E(sy,sx)
 =-(-1)^{(p-1)(q-1)+q+pq}\tau_E(xy)
 =(-1)^p\tau_E(xy).
\]
A nonzero value has $p+q=d$, so the two tangent degrees sum to $d-2$.
This proves both the alternating sign and the displayed target shift.

For chain compatibility, use $d(sx)=-s(dx)$ and $\tau_E(d(xy))=0$.
The chain-map equation becomes $(-1)^p\tau_E(d(xy))=0$.
The underlying trace pairing is perfect, and shifting preserves its
nondegeneracy.
\end{proof}

The closedness of the global form has more content than this linear
calculation. To specify it, let
\[
 \operatorname{DR}(\mathcal M)
 =\operatorname{R\Gamma}\!\left(
 \mathcal M,\prod_{q\geq0}(\textstyle\bigwedge^q\mathbb L_{\mathcal M})[-q]
 \right)
\]
be the Hodge-completed derived de Rham complex. On its $q$th factor
the internal differential is $(-1)^q d$. The de Rham differential
goes to the next factor and has total degree one.
Its filtration $F^p$ retains the terms with $q\geq p$.

A closed $p$-form of degree $n$ is a point of the derived mapping space
$\operatorname{Map}(k[-n],F^p\operatorname{DR}(\mathcal M)[p])$.
Its cocycle has compatible components in every Hodge weight at least $p$.
Projection to the first term gives the underlying $p$-form. A closed
two-form is \emph{$n$-shifted symplectic} if its underlying pairing
induces an equivalence
$\mathbb T_{\mathcal M}\simeq\mathbb L_{\mathcal M}[n]$.

\begin{theorem}[Brav--Dyckerhoff]\label{cygeom:thm:moduli-form}
Let $\mathcal C$ be of finite type and equipped with a left
Calabi--Yau structure of dimension $d$. Its moduli of pseudo-perfect
objects carries a $(2-d)$-shifted symplectic form. With the suspension
convention above, its underlying pairing at $E$ is
\eqref{cygeom:eq:suspended-pairing}.
\end{theorem}

\begin{proof}[Construction and nondegeneracy]
The universal family induces a circle-equivariant map on Hochschild
chains. Pass to negative cyclic complexes and take the weight-two
component of the derived de Rham realization. This sends $\eta$ to a
closed two-form of degree $2-d$, by
\cite{cygeom:BD-moduli}*{Proposition~5.2}. Its higher components
retain the cyclic coherence in $\eta$.

The underlying form is the
composition of endomorphisms followed by the induced trace
\cite{cygeom:BD-moduli}*{Proposition~5.3}. Under
\eqref{cygeom:eq:tangent}, suspension gives the sign in
\eqref{cygeom:eq:suspended-pairing}. The bimodule nondegeneracy of
$\eta$ makes this trace pairing perfect, hence identifies the tangent
complex with the shifted cotangent complex. This is
\cite{cygeom:BD-moduli}*{Theorem~5.5(1)}.
\end{proof}

For the origin in affine three-space, the tangent cohomology is therefore
\[
 H^{-1}(\mathbb T_E)=k,\quad
 H^0(\mathbb T_E)=k^3,\quad
 H^1(\mathbb T_E)=\bigwedge^2 k^3,\quad
 H^2(\mathbb T_E)=k.
\]
The form has degree $-1$. It pairs the first and last spaces and the
middle two, with the explicit signs above. The ordinary family of points
is smooth, but this does not remove the higher tangent groups of the
derived object moduli. The same formula gives degree zero for a
Calabi--Yau surface and degree two for a zero-dimensional Calabi--Yau
category.

\section{Proper support and compact geometry}
\label{cygeom:sec:support}

Over $A=k[x_1,x_2,x_3]$, the free module has infinite-dimensional
endomorphisms and is not a pseudo-perfect $k$-point. A module with
finite-dimensional total cohomology is pseudo-perfect. Such a module is
also perfect over $A$. Tensor a finite diagonal
resolution with the module. Each resulting free $A$-term has a
perfect $k$-complex of coefficients. The point calculation thus belongs
to a proper-support subcategory of a smooth nonproper category.

For comparison, let $X$ be a smooth projective variety of dimension
$d$ with a chosen trivialization of its canonical line $\omega_X$.
Serre duality gives, for perfect complexes $E,F$,
\begin{equation}\label{cygeom:eq:serre}
 \operatorname{Ext}^i_X(E,F)^*
   \simeq\operatorname{Ext}^{d-i}_X(F,E).
\end{equation}
Here the usual functor $(-)\otimes\omega_X[d]$ becomes the shift
$[d]$ by the chosen trivialization. This trace is a geometric input.
When a closed form on moduli is used, its compatible cyclic structure
is retained as in Theorem~\ref{cygeom:thm:moduli-form}.

\begin{corollary}\label{cygeom:cor:no-tilting}
In a nonzero proper $d$-Calabi--Yau category with $d>0$, no nonzero
object $T$ has $\operatorname{Ext}^i(T,T)=0$ for every $i\ne0$.
In particular, there is no tilting generator with this vanishing.
\end{corollary}

\begin{proof}
The identity makes $\operatorname{Hom}(T,T)$ nonzero. Duality identifies
its dual with $\operatorname{Ext}^d(T,T)$, which is therefore nonzero.
\end{proof}

A compact generator can still exist. Its endomorphisms must retain their
dg structure, including the class dual to its identity. Local tilting
presentations on nonproper spaces have different hypotheses and do not
contradict this corollary.

\section{Determinant lines along extensions}
\label{cygeom:sec:orientation}

For a perfect complex $V$, its determinant is
$\det V=\bigotimes_i(\det H^iV)^{(-1)^i}$. An exponent $-1$ denotes
the dual line. We use determinant superlines: their parity is
$\chi(V)\bmod2$, and interchanging lines of parities $a,b$ has sign
$(-1)^{ab}$. The determinant functor is multiplicative on exact
triangles, with these symmetry signs.

On the object moduli of a three-dimensional Calabi--Yau category set
$\mathcal D_E=\det\operatorname{RHom}(E,E)$. Since
$\mathbb T_E=V_E[1]$, the determinant of the cotangent complex is
$\det\mathbb L_E=\det(V_E^*[-1])=\mathcal D_E$.
If $\operatorname{Ext}^i(E,E)=0$ outside $0\leq i\leq3$, duality gives
\begin{equation}\label{cygeom:eq:det-square}
 \mathcal D_E\simeq
 \left(\det\operatorname{Ext}^0(E,E)\otimes
       (\det\operatorname{Ext}^1(E,E))^{-1}\right)^{\otimes2}.
\end{equation}
Indeed, the degree-two space is dual to degree one and degree three
is dual to degree zero. The parity of $V_E$ is even, as the same
pairings show.

Extensions impose a further relation. Consider an exact triangle
$E_1\to E\to E_2\to E_1[1]$ on which all Hom complexes are perfect
over $k$. Put $U=\operatorname{RHom}(E_1,E_2)$.
Applying Hom in each variable and the determinant functor gives
\[
 \mathcal D_E\simeq
 \mathcal D_{E_1}\otimes\mathcal D_{E_2}\otimes
 \det U\otimes\det\operatorname{RHom}(E_2,E_1).
\]
The factors are ordered as displayed using the determinant symmetry.

Calabi--Yau duality identifies the last complex with $U^*[-3]$.
Dualizing and shifting by an odd integer each invert a determinant.
Together they leave it unchanged. Consequently
\begin{equation}\label{cygeom:eq:triangle-determinant}
 \mathcal D_E\simeq
 \mathcal D_{E_1}\otimes\mathcal D_{E_2}\otimes(\det U)^{\otimes2}.
\end{equation}

An \emph{orientation datum} for these extension spaces consists of a
superline bundle $\mathcal L$ with $\mathcal L^{\otimes2}\simeq\mathcal D$
and isomorphisms
\begin{equation}\label{cygeom:eq:orientation}
 \mathcal L_E\simeq
 \mathcal L_{E_1}\otimes\mathcal L_{E_2}\otimes
 \det\operatorname{RHom}(E_1,E_2)
\end{equation}
whose squares are \eqref{cygeom:eq:triangle-determinant}. These are the
square-root and triangle relations of
\cite{cygeom:KS-orientation}*{\S5.2}.

We require the choices to be compatible with base change and with
successive extensions. For a flag with three successive quotients,
\eqref{cygeom:eq:orientation} has two iterated composites.
These must coincide, using the determinant symmetry just fixed.
The analogous coherences are retained
on longer flags.

Formula~\eqref{cygeom:eq:det-square} supplies a square root on a stratum
where its cohomology spaces form vector bundles. It does not supply the
gluing across dimension jumps or the flag isomorphisms
\eqref{cygeom:eq:orientation}. Those are part of the orientation datum.
Once such a datum is chosen, a Hall construction can use the resulting
oriented coefficients on extension correspondences. The pullback,
direct image and comparison of these coefficients must still be
constructed for the chosen cohomology theory. These maps live on flag
spaces and their incidence correspondences.
